\section*{الفصل \ref{ch:b2:limb-organogenesis} --- تكوّن الأعضاء: بناء طرف رباعي الأرجل}

\begin{solution}{exo:b2:limb-organogenesis:1}
الداني--القاصي: \omterm{def:b2:limb-organogenesis:bud}{العرف الأديمي القمي} على امتداد القمة، وعوامل FGF.
والأمامي--الخلفي: \omterm{def:b2:limb-organogenesis:bud}{منطقة النشاط المستقطب} عند الحافة الخلفية،
\omterm{prop:b2:limb-organogenesis:zpa}{وSonic hedgehog}. والظهري--البطني: الأديم الظاهر الظهري، وWnt7a.
\end{solution}

\begin{solution}{exo:b2:limb-organogenesis:2}
القطعة الجذعية: العضد، والفخذ (Hoxa/d 9--10). والقطعة الوسطى: الكعبرة والزند،
والظنبوب والشظية (Hox 11). والقطعة الطرفية: الرسغ واليد، والكاحل والقدم
(Hox 13).
\end{solution}

\begin{solution}{exo:b2:limb-organogenesis:3}
تحديد حقل الطرف (Tbx5)؛ ثم نمو البرعم خارجًا تحت عامل FGF العرفي؛
ثم تنميطه \omterm{def:b2:limb-organogenesis:bud}{بمنطقة النشاط المستقطب} والأديم الظاهر الظهري؛ ثم تكثف الخلايا
المغادرة \omterm{prop:b2:limb-organogenesis:aer}{لمنطقة التقدم} أشعةً؛ ثم تمايز الغضروف؛ ثم تكوّن المفاصل
حيث يُكبت الغضروف؛ ثم غزو الأوعية الدموية وحلول العظم
محل الغضروف؛ ثم موت الخلايا بين الإصبعية؛ ثم هجرة طلائع العضل من
القُطَيعات واندماجها وارتباطها بأوتار.
\end{solution}

\begin{solution}{exo:b2:limb-organogenesis:4}
(أ) جذمور بعضد ضئيل أو بلا عضد؛ (ب) عضد وكعبرة وزند
بلا يد؛ (ج) طرف كامل --- إذ يحل FGF محل العرف.
\end{solution}

\begin{solution}{exo:b2:limb-organogenesis:5}
$x = 40\ln 2 = \qty{28}{\micro m}$، و$40\ln 5 = \qty{64}{\micro m}$،
و$40\ln 20 = \qty{120}{\micro m}$: أي أقاليم من 28 و36 و
\qty{56}{\micro m}.
\end{solution}

\begin{solution}{exo:b2:limb-organogenesis:6}
(أ) 4-3-2-2-3-4، أي طقم مرآوي كامل. (ب) المنبع الأضعف يعطي
ذروة أدنى، فلا يُضاف إلا الإصبع ذو العتبة المنخفضة: أي 2-2-3-4. (ج)
والتعرض القصير يحدد كذلك الإصبع الزائد الأمامي وحده:
أي 2-2-3-4 --- فالخلايا تكامل التركيز على الزمن.
\end{solution}

\begin{solution}{exo:b2:limb-organogenesis:7}
$5000\times 2^{8} = 1.3\times 10^{6}$. وبلوغ $4\times 10^{6}$ يحتاج إلى
$\log_2 800 = 9.6$ مضاعفة، فلا بد أن يكون متوسط الدورة نحو
\qty{10}{h}، أو أن يشمل العدّ طلائع العضل التي
تهاجر إلى الداخل من القُطَيعات؛ أما موت الخلايا فيزيد التباين
سوءًا لا إصلاحًا.
\end{solution}

\begin{solution}{exo:b2:limb-organogenesis:8}
خلايا الكتكوت بين الإصبعية تتلقى إشارة BMP فتموت؛ وخلايا البطة
بين الإصبعية تعبّر عن مثبط BMP فتنجو، فيبقى الغشاء.
والمثبط على حبة ينقذ غشاء الكتكوت: فتكون القدم
مكفَّفة. والموت في موعده في طُعم يعني أن الخلايا كانت
قد عُلّمت قبل النزع --- فالقرار يُتخذ باكرًا
ويُنفَّذ لاحقًا، ذاتيًا.
\end{solution}

\begin{solution}{exo:b2:limb-organogenesis:9}
بلا Hoxa13 وHoxd13: لا قطعة طرفية --- فينتهي الطرف الأمامي عند
الرسغ. وبلا Hoxa11 وHoxd11: كعبرة وزند قصيران جدًا
ويد ملتصقة بالعضد مباشرة تقريبًا. فكل مقطع
يحتاج زوج Hox الخاص به؛ والنطق ليست مجرد علامات بل
تعليمات.
\end{solution}

\begin{solution}{exo:b2:limb-organogenesis:10}
عامل FGF من العرف يبقي Shh مشتعلًا في \omterm{def:b2:limb-organogenesis:bud}{منطقة النشاط المستقطب}؛ وShh، عبر مرحّل
في الأديم المتوسط، يبقي FGF مشتعلًا في العرف. فإشارة كل مركز
تشهد أن الآخر يعمل، فلا ينمو طرف قط دون
أن يُنمَّط ولا يُنمَّط دون أن ينمو، ويصلح كل منهما فقدًا
عابرًا للآخر. ولو عجز Shh عن صون
FGF لضمر العرف في غضون يوم ولتوقف النمو الخارجي
عند أي مقطع بلغه: أي طرف مبتور. أما حلقة
المخاض فينهيها الحدث الذي تقوده --- إذ تزيل الولادة
المنبه؛ وحلقة الطرف ينهيها التمايز --- إذ يكف
الأديم المتوسط عن الترحيل حين يتم الطرف.
\end{solution}

\begin{solution}{exo:b2:limb-organogenesis:11}
$e^{4r} = 8$: ومنه $r = \ln 8/4 = \qty{0.52}{d^{-1}}$. أما النمو:
فتكاثر أطول وأسرع لغضروف الإصبع (بمعزِّز طرفي
يرفع تعبير مورّثة نمو في الأصابع). وأما الموت: فنجاة
النسيج بين الإصبعي بالتعبير عن مثبط BMP،
وهو ما يبقي الغشاء.
\end{solution}

\begin{solution}{exo:b2:limb-organogenesis:12}
تُعاد الإشارات لأن وحدة مستقبل--نقل عاملة
رخيصة إعادة النشر: فالذي يتغير بين الأنسجة ليس
الإشارة بل طائفة المورّثات التي تكون الخلايا المستقبِلة أهلًا
لإشعالها، فيعني منحى التركيز نفسه عصبونًا محركًا في
موضع والإصبع 4 في آخر. ولذلك تصيب \omterm{def:b2:genome-diversification:mutation}{الطفرة} في مورّثة كهذه
أعضاء كثيرة معًا --- فطفرات Shh تعطي
الدماغ الأمامي الشمولي وعيوب الأطراف معًا --- إلا إذا وقعت في
معزِّز نوعي النسيج، فتغيّر عندئذ عضوًا واحدًا فقط،
وهكذا يقع معظم التغير التطوري في هذه المورّثات.
\end{solution}

\begin{solution}{pb:b2:limb-organogenesis:1}
\textbf{1.} $96/12 = 8$ مضاعفات: ومنه $2000\times 256 = 5.1\times 10^{5}$
خلية.
\textbf{2.} $1000/256 \approx 4$: أي تضخم الخلايا وحشوة الغضروف.
\textbf{3.} الطول $\times 10$؛ وأسطوانة ثابتة نصف القطر كانت
لتعطي $\times 10$ في الحجم؛ ومئة الضعف الزائدة تعني أن نصف القطر
نما كذلك نحو عشرة أضعاف --- فقد اتسع البرعم مجدافًا وهو
يطول.
\textbf{4.} عند 3 أيام $300/400 = \qty{75}{\%}$؛ وعند 7 أيام $300/4000 =
\qty{7.5}{\%}$.
\textbf{5.} القمة وحدها تحت الإشارة؛ ومع تقدم القمة،
تغادر الخلايا المتروكة خلفها المنطقةَ بالترتيب الذي صُنعت به، فتُحدد
البنى الدانية أولًا والقاصية آخرًا.
\textbf{6.} القطعة الوسطى (الكعبرة والزند)، المحددة بين 3.5 و
4 أيام.
\textbf{7.} القطعة الجذعية بحلول 3.5 أيام، والوسطى بحلول 4.0 أيام، والطرفية بحلول 4.5
أيام (وأطراف الأصابع بعد ذلك).
\textbf{8.} نصف يوم: أي دورة خلوية واحدة.
\textbf{9.} كل دورة تُقضى في \omterm{prop:b2:limb-organogenesis:aer}{منطقة التقدم} تقدّم شفرة
Hox (9--10، ثم 11، ثم 13)؛ والخلايا المغادرة بعد دورة أو دورتين أو
ثلاث تحمل شفرة القطعة الجذعية أو الوسطى أو الطرفية.
\textbf{10.} جناح كامل: فعامل FGF هو ما كان يوفره العرف.
\textbf{11.} نمو خارجي ثانٍ من السطح الظهري: أي بنى قاصية
مضاعفة، وطرف زائد.
\textbf{12.} يكف الأديم المتوسط عن ترحيل Shh إلى العرف، فينتهي التعبير
عن Shh، ويضمر العرف، وتتمايز الخلايا
بدل أن تنقسم: فتُكسر الحلقة التي كانت تصون النمو.
\textbf{13.} $k = \ln 2/1740 = \qty{4.0e-4}{s^{-1}}$؛ و$\lambda =
\sqrt{10^{-12}/4\times 10^{-4}} = \qty{50}{\micro m}$.
\textbf{14.} $x = 50\ln(1/0.6) = 26$، و$50\ln 4 = 69$، و$50\ln 12.5 =
\qty{126}{\micro m}$: أي أقاليم من 26 و43 و\qty{57}{\micro m}.
\textbf{15.} $600 - 126 = \qty{474}{\micro m}$.
\textbf{16.} يتحرك كل حد إلى الخارج بمقدار $50\ln 2 = \qty{35}{\micro
m}$: فيصير 61 و104 و161. والعروض الثلاثة 61 و43 و\qty{57}{\micro m}:
فلا ينمو إلا إقليم الإصبع 4، وتُزاح البقية إزاحةً فحسب.
فينزاح النميط أماميًا بدل أن يكسب إصبعًا، وتنكمش
المنطقة الخالية من الأصابع من 474 إلى \qty{439}{\micro m}.
\textbf{17.} يحتاج كل طقم إلى \qty{126}{\micro m}: أي
\qty{252}{\micro m} على الأقل؛ وفي برعم \qty{600}{\micro m} متسع، مع
\qty{348}{\micro m} من النسيج الخالي من الأصابع في الوسط.
\textbf{18.} عند \qty{200}{\micro m} يتداخل التدرجان و
يتجمعان: فترى الخلايا الوسطى ما يكفي للإصبع 3 وتزول أقاليم
الإصبع 2 --- أي 4-3-4 أو 4-3-3-4، وهي أصابع أقل من طقم
مرآوي كامل.
\textbf{19.} ربع الخلايا يعطي $c_0/4$: فلا تبلغ أي خلية
عتبة الإصبع 4 ولا الإصبع 3 (أي $0.25 c_0$ عند الحافة نفسها)، و
تُتجاوز عتبة الإصبع 2 حتى $50\ln(0.25/0.08) =
\qty{57}{\micro m}$: أي إصبع 2 زائد واحد.
\textbf{20.} $3\times 300\times 300\times 200 = 5.4\times 10^{7}$
\unit{\micro m^3}: أي $5.4\times 10^{4}$ خلية؛ وعلى مدى \qty{86400}{s}،
$0.6$ خلية في الثانية.
\textbf{21.} حبة BMP في غشاء البطة: يموت الغشاء وتنفصل
الأصابع؛ والمثبط في غشاء الكتكوت: ينجو الغشاء، فتكون القدم
مكفَّفة.
\textbf{22.} $e^{0.5\times 4} = e^{2} = 7.4$. والتغير الثاني
نجاة الغشاء بين الإصبعي بمثبط BMP.
\textbf{23.} الطور المتأخر من التعبير عن Hox 13 الذي يحدد
قطعة طرفية، وبقاء العرف: فطية السمكة تصنع أشعة
زعنفية بدلًا من ذلك. وزعنفة \emph{Tiktaalik} تحوي عظامًا شبيهة بالرسغ --- أي
زعنفة تبدأ أن تصير طرفًا.
\textbf{24.} $5.4\times 10^{4}/1000 = 54$: ومنه $\log_2 54 = 5.8$ مضاعفة،
أي نحو \qty{69}{h}، أي ثلاثة أيام لصنع ما يزيله يوم واحد.
\textbf{25.} ثماني مضاعفات؛ وحدود الأصابع عند 26 و69 و
\qty{126}{\micro m}؛ والتضاعف المرآوي يحتاج إلى
\qty{252}{\micro m} على الأقل.
\end{solution}
